\chapter*{Úvod} % chapter* je necislovana kapitola
\addcontentsline{toc}{chapter}{Úvod} % rucne pridanie do obsahu
\markboth{Úvod}{Úvod} % vyriesenie hlaviciek

Posledných niekoľko desaťročí prebieha informatizácia spoločnosti, v súčasnosti je vo veľmi pokročilom stave a drvivé
množstvo informácií (takmer 100\%) je zaznamenávaných a spracovávaných v digitálnej forme. Potenciál digitálnych
informačných a komunikačných technológií(d-IKT) si uvedomuje aj Európska únia(EÚ). Rozvoj znalostnej spoločnosti s trvalo
udržateľným rozvojom sú hlavnými prioritami EÚ, ktoré sú rozpracované vo viacerých strategických dokumentoch, pre 
informatizáciu spoločnosti sú najdôležitejšie \enquote{Digital Single Market Strategy}\footnote{Stratégia cielená na
voľný pohyb digitálneho tovaru a odstraňovania prekážok v rámci Európy.} a \enquote{Digital Agenda for Europe}\footnote{Agenda
politík určených na maximalizovanie prínosu digitálnych technológií.}. 

EÚ si je vedomá, že d-IKT sa stali prvkami kritickej infraštruktúry spoločnosti, je potrebné ich chrániť a preto prijala
\enquote{European Cybersecurity Strategy}\footnote{Stratégia cielená na budovanie odolnosti voči kybernetickým hrozbám.},
\enquote{Network and Information Security Directive (NIS Directive)}\footnote{Smernica o bezpečnostných opatreniach
pre zabezpečenie vysokej úrovene informačnej bezpečnosti.} a \enquote{EU Cybersecurity Act:}\footnote{Zákon na posilnenie
agentúry ENISA a vytvorenie certifikačného rámca pre digitálne produkty.}.

Ani Slovensko nezaostalo a tiež sa snaží o informatizáciu spoločnosti cez dokumenty ako \enquote{Stratégia digitálnej
transformácie Slovenska 2030} a \enquote{Akčných plánov digitálnej transformácie Slovenska}, pre oblasť kybernetickej
a informačnej bezpečnosti to sú \enquote{Národný program pre ochranu a obranu kritickej infraštruktúry} a
\enquote{Národná stratégia a akčný plán kybernetickej bezpečnosti}.

Nariadenia EÚ spolu s národnými stratégiami, koncepciami a akčnými plánmi sa premietli do viacerých zákonov Slovenskej
republiky(SR). Tieto zákony podrobnejšie definujú pravidlá a upravujú povinnosti subjektov v kybernetickom priestore
Slovenska a EÚ, sú to napr. 45/2011 Z. z.(zákon o kritickej infraštruktúre), 272/2016 Z. z.(zákon o elektronickom podpise)
alebo 18/2018 Z. z.(zákon o ochrane osobných údajov).

Najdôležitejšie zákony a vykonávacie predpisy SR stanovujúce povinnosti subjektov v kybernetickom priestore sú:
\begin{itemize}
  \item 69/2018 Z. z., zákon o kybernetickej bezpečnosti,
  \item 362/2018 Z. z., vyhláška Národného bezpečnostného úradu, ktorou sa ustanovuje obsah bezpečnostných opatrení,
    obsah a štruktúra bezpečnostnej dokumentácie a rozsah všeobecných bezpečnostných opatrení,
  \item 95/2019 Z. z., zákon o informačných technológiách vo verejnej správe a
  \item 179/2020 Z. z., vyhláška Úradu podpredsedu vlády Slovenskej republiky pre investície a informatizáciu, ktorou
    sa ustanovuje spôsob kategorizácie a obsah bezpečnostných opatrení informačných technológií verejnej správy.
\end{itemize}

Vyššie uvedené právne normy určujú povinnosti v kybernetickom priestore veľkému počtu vlastníkov, prevádzkovateľov a
správcov informačných systémov. Národný bezpečnostný úrad (podľa kompetenčného zákona 69/2018 Z.z. \cite[\S 5, ods. 1]{69/2018})
zodpovedá za kybernetickú bezpečnosť a nedávno zaradil do registra prevádzkovateľov základných služieb mestá, obce,
mestské časti a niektoré služby a organizácie spadajúce pod mestá. To znamená, že okrem informačných systémov verejnej
správy(ISVS) ústredných štátnych orgánov sa prvkami kritickej infraštruktúry (podľa 69/2018 Z.z) stali aj ISVS malých 
miest, obcí a organizácií, ktoré spadajú do najnižšej kategórie I (podľa 69/2018 Z.z a podľa 95/2019 Z.z.).


Tu sa dostávame k veľkému problému pre SR, ktorý vyplýva z tohto zákona. Jednou z minimálnych požiadaviek 69/2018 Z.z
\cite[\S 20, ods. 4, písm a)]{69/2018} je ustanovenie manažéra kybernetickej bezpečnosti. To znamená, že všetky subjekty
zaradené do zoznamu prevádzkovateľov základnej služby potrebujú odborne kvalifikovaného zamestnanca zodpovedného za ochranu
informačných systémov a dodržiavania legislatívnych povinností. Vyjadrené číselne, ide zhruba o 20 000 ľudí, ktorých SR
nemá k dispozícií a v blízkej budúcnosti ani nebude mať\footnote{Pre tak veľké množstvo odborníkov nestačia kapacity
vzdelávacích inštitúcií, FMFI UK pripraví ročne menej ako 10 absolventov.}. Tým sme sa dostali
k cieľom bakalárskej práce:
\begin{itemize}
  \item analyzovať legislatívne požiadavky na ochranu ISVS,
  \item identifikovať povinnosti prevádzkovateľa ISVS kategórie I, 
  \item navrhnúť postup, ako splniť tieto povinnosti a
  \item vytvoriť pomocný systém pre manažéra KIB, ktorý mu pomôže si splniť legislatívne povinnosti.
\end{itemize}
V nasledujúcich kapitolách sa pokúsime tieto ciele splniť a na konci zhodnotíme, do akej miery sa nám aj naozaj podarilo
ich splniť. Najprv začneme s terminológiou, aby sme začali a následne budovali na pevnom základe. 

%Bakalárskej práca sa skladá z kapitol:
%\begin{itemize}
%  \item základné pojmy, kde oboznámime čitateľa so základnými pojmami informačnej bezpečnosti,
%  \item analýza právnych požiadaviek, kde analyzujeme vyššie uvedené zákony a vyhlášky kybernetického priestoru SR,
%  \item zavádzania ISMS, kde popíšeme kroky potrebné na zavedenie systému riadenia informačnej bezpečnosti,
%  \item systém, kde popíšeme navrhnutý systém a jeho jednotlivé časti,
%  \item nedostatky a kam ďalej, kde rozoberieme  
%\end{itemize}
